\section{The Succubus and the Houri}

\paragraph{Princeton Gnosis}

The French philosopher \textbf{Raymond Ruyer} wrote a strange book, the \emph{Gnosis of Princeton}, that intrigued thinkers such as \textbf{Henry Corbin} and even Jean Borella\footnote{\url{http://jeanborella.blogspot.com/2008/11/la-gnose-ruyrienne-religion-de-lge.html}}. It is a work of philosophical fiction disguised as fact, that left more than a few convinced that he was describing an actual movement. In any case, it establishes fiction as a legitimate philosophical genre.

The theme is that an elite cadre of scientists from Princeton, having rejected any positivist and materialist interpretations of science, set out to create a neo-gnostic religion based on Ruyer's peculiar philosophical positions, e.g., pansychism, pantheism, holons, and so on.

This raises an intriguing question. Suppose you encountered a movement dedicated to esoteric gnosis, how would you --- in your current state of development --- be able to distinguish the ``real thing'' from a hoax?

\paragraph{Spirituality for the Last Man}


In a book dictated by \textbf{Orin}, an alleged ``source of wisdom'', we learn that 
\begin{quotex}
loving yourself means accepting yourself as you are right now ... It is an agreement with yourself to appreciate, validate, accept and support who you are at this moment.
\end{quotex}


There is no exhortation to self-knowledge, to striving, to overcoming and transcending one's spontaneous instinctual and emotional arisings, to developing a true will. As you are ``right now'' is little more than the vector sum of unknown forces encountered throughout the course of life. A real source of wisdom would show how to develop a Real I that can be the master of such forces.

\paragraph{Facing the Facts}


An old Latin saying is, ``Inter faeces et urinam nascimur''. This applies also to our spiritual life. If we accept living spiritually in our own mess, we will never be motivated to clean up. That is not negative thinking but leads to an authentic Joy. This was channeled to me by a source much higher than Orin.

\paragraph{Attack of the Succubus}


A frustrating aspect of spiritual development is backsliding, i.e., reverting to the same unhelpful states of consciousness repeatedly. I can go for days with a mind as still as the ocean on a calm summer's day, interspersed with periods of bliss. Yet, unexpectedly, a demonic attack occurs, usually in that twilight state between waking and sleeping.

In a way, I am fortunate to be so weak, because I could not withstand the demonic attacks borne, for example, by a Saint Anthony\footnote{\url{http://www.newadvent.org/fathers/2811.htm}} or an Anne Catherine Emmerich\footnote{\url{http://www.spiritdaily.net/amazingemmerich.htm}}.

The Succubus attacks by bringing sexual imagery into consciousness. Now it doesn't help to avoid pornography, which in any case does not interest me. Rather, she uses events from my own memory. It's odd how random encounters from decades ago can be so vivid, while the memories of more important life events have dimmed. Now it is the case that sex creates one flesh, which is the problem with promiscuity. Even if you haven't seen a former sexual partner in several years, it is often very easy to slip back into that relation as though just a month passed by.

In the \emph{Life of Saint Anthony}, we never do find out exactly how the demonic attack was resisted. Perhaps, in this case, it may be useful to reveal the method. First of all, it is necessary to become aware that the pleasure of the imagery is false. Bring to mind the decay of the flesh, bodily excretions, and so on, in as much detail as you need.

Guenon claims that the personal element is irrelevant. That would make sense except that we necessarily begin in the human state. So some concrete suggestions may help others, or probably most will just be offended. So I hesitated to post such personal information, until I came to the realization that others have encountered the same Succubus. A Houri is the opposite, a female figure who leads upwards to Paradise. A Hadith claims that the Houri does not menstruate, urinate, or defecate. It is almost as though a Houri were interceding for me.

\paragraph{Life in Hell}


To the modern mind, the idea of being condemned to Hell seems unjust, something only a cruel God could do. Yet, what seems to be a conundrum for exoteric thought is understood differently esoterically. Hell is a postmortem state that is entirely just. Is a hydrometer unfair for indicating the alcohol content of your beer?

Actually, the time element is irrelevant since it confuses eternity with unending time. If a punishment is itself cruel, then why would ending it after a million years be an act of mercy? A thousand years? A day? Just once is already cruel.

Note that descriptions of the Hell realms described, for example, by Dante, are repeated exactly the same over and over. Since repetition is impossible for the Infinite, we should interpret that as a static state of being. A dead body can no longer act, improve or evolve. Likewise, a spiritually dead soul cannot act, improve, or evolve. It is stuck at the same state, at least until the passing away of this world.

The soul experiencing the Beatific Vision, on the other hand, is alive. The experience is never duplicated, since God is Infinite.

\paragraph{Liberalism and Strange Attractors}


In a dynamical system, an attractor\footnote{\url{https://en.wikipedia.org/wiki/Attractor}} is a point toward which the system tends to evolve for a wide variety of starting conditions. In a fallen world, Liberalism, or the Consensus Reality of the Modern World, is such an attractor. On the one hand, atheists, secularists, and materialists tend to end up as Liberals. Curiously, New Agers and liberal Christians, despite starting from a widely different set of assumptions from the secularists, are likewise attracted to Liberalism.

In this case, the attractor is like a steep bank that leads into the sea.

\paragraph{True Man and Social Justice Warrior}


When I was at St Patrick's church in my youth, a young priest named Fr Shanley arrived. The name will be familiar to you if you've seen the recent movie \emph{Spotlight} about clerical abuse in the Boston area. They always speak of ``children'' but the truth is that most cases involved pederasty, i.e., sexual activity with an adolescent under the age of consent. I just looked it up on Wikipedia and was surprised to learn that pederasty is

\begin{quotex}
not likely to give rise to psychological discomfort or neuroses for all or most males.
\end{quotex}


Just so you know what's coming. If that be true, then the large damage awards make no sense.

Well, back to the point. First of all, I had no interactions with him other than at mass or at CCD classes. Nor was he involved in with any of the youths in my circle, probably because we were too old. He hung out with the 10th graders. I suspect he was able to identify SSA boys, whom he chose as his targets, through their confessions. Otherwise, I don't understand why they capitulated so easily or why he was never severely beaten up. That was the customary way of dealing with unwanted sexual advances in those days.

He was the new breed of priest, youthful, long hair, jeans wearing, with ``social concerns''. He would be called today a Social Justice Warrior. That designation has a certain élan today, since it costs nothing personally, just a willingness to attack others. That gave him the cover for his misdeeds; he was too cool to be bad, or at least he was ``good bad''.

In the Tao Te Ching\footnote{\url{https://www.gornahoor.net/?s=tao\%20te\%20ching}}, the \textbf{True Man} is not humane, benevolent or compassionate. From the Taoist perspective, I suppose, Fr. Shanley was ``benevolent'' and ``humane'', or at least gave the appearance of it. Those are the types to be wary of.

That is the opposite of the True Man, who has overcome such sentimentality.

\paragraph{Clash of Religion Traditions}


From reading Michael McClain, it is unclear if he is promoting Persian Sufism, Russian Orthodoxy or old-fashioned Spanish-style Catholicism. When I asked him which Tradition he followed, he refused to answer, although it is easy to guess. I found this in an old letter from him to a British Islamic publication\footnote{\url{http://aaiil.org/text/articles/thelight/1983/light19830408.pdf}}.

\begin{quotex}
It is Islam which is best equipped to save what can be saved from this possible wreck, and for this reason every traditional Christian should do everything in his power to favour the strong presence of Islam in Europe and America. (April 1983)
\end{quotex}


Obviously, the situation is quite different now from what it was 33 years ago. Now at the level of depth, two esoterists of two different traditions can understand each other, and even get along. However, at the exoteric level, whenever two Traditions clash, one would have to be the master of the other.

Perhaps he was being true to the ideal that the important thing is the preservation of Tradition rather than any particular form.

God knows best.


\flrightit{Posted on 2016-08-15 by Cologero}

\begin{center}* * *\end{center}

\begin{footnotesize}
\begin{sffamily}
\texttt{Wayne Ferguson on 2016-08-15 at 07:21 said:}

``Inter faeces et urinam nascimur''. St Augustine of Hippo

Trying to verify the source--does anyone happen to have a citation?

\url{https://en.wikiquote.org/wiki/Talk:Bernard\_of\_Clairvaux}


\url{https://it.wikipedia.org/wiki/Inter\_faeces\_et\_urinam\_nascimur}


\hfill

\texttt{Cologero on 2016-08-15 at 14:21 said:}

Wayne, probably nobody ever said it. But is it true?

\hfill

\texttt{Matt on 2016-08-21 at 03:12 said:}

Ruyer's work seems quite intriguing. If there are any english translations of his works I'll have to check them out at some point. He seems to be a rarity among the modern French intelligenstia in that he argued in favour of a teleological worldview, which is not in line with the fashionable existentialism and post-modernism that his peers flocked towards.

\hfill

\texttt{Max on 2016-09-04 at 11:21 said:}

There seems to be a collective element to the attack of that sort of imagery. That would imply more far reaching consequences than if the phenomena were just isolated to us, and gives a deepened sense of responsibility to conduct the spiritual battle.

\hfill

\texttt{Ulysses on 2016-09-24 at 20:22 said:}

A French hadith says that the houris supplicate: « O Seigneur, aide le dans la voie de ta religion, fais que son cœur t'obéisse et qu'il parvienne jusqu'à nous. Ô Toi le plus Miséricordieux des Miséricordieux » ... . still looking for a precise reference ... 

\hfill

\end{sffamily}
\end{footnotesize}

